\chapter{SIMPULAN DAN SARAN}

\section{Simpulan}
Berdasarkan hasil penelitian mengenai integrasi \textit{Retrieval-Augmented Generation} (RAG) dan deteksi objek multimodal pada sistem agentik untuk identifikasi penyakit tanaman, dapat ditarik beberapa simpulan sebagai berikut:

\begin{enumerate}[nosep]
    \item Integrasi model deteksi objek dan pencarian multimodal berhasil diimplementasikan dalam arsitektur sistem agentik. Model YOLOv11-\textit{small} terbukti menjadi varian paling optimal untuk deteksi daun (\textit{closed-set}) dengan menyeimbangkan akurasi mAP@50 sebesar \qty{89.16}{\percent} dan latensi inferensi sekitar \qty{47}{ms}. Sementara itu, integrasi model OWLv2 memberikan fleksibilitas deteksi \textit{open-vocabulary} untuk objek non-daun, meskipun dengan konsekuensi latensi yang lebih tinggi. Penggunaan model SCOLD pada modul pencarian gambar-teks memungkinkan identifikasi penyakit berbasis kemiripan visual dan semantik yang efektif.
    \item Mekanisme \textit{Retrieval-Augmented Generation} (RAG) berhasil diterapkan untuk memperkaya konteks jawaban agen dengan informasi faktual. Konfigurasi pencarian hibrida yang menggabungkan BM25 dan \textit{embedding} Gemini, serta diperkuat oleh \textit{reranker} Voyage AI, menghasilkan kinerja pengambilan informasi terbaik dengan skor nDCG@5 mencapai 0,954. Mekanisme ini memastikan agen memiliki landasan pengetahuan yang relevan sebelum memberikan diagnosis atau saran penanganan.
    \item Perancangan agen \textit{Large Language Model} (LLM) menggunakan pola \textit{Reasoning and Acting} (ReAct) pada kerangka kerja LangGraph terbukti mampu mengoordinasikan penggunaan alat secara adaptif. Penerapan instruksi sistem berbasis CO-STAR dan aturan kebenaran (\textit{truthfulness rules}) berhasil menjaga perilaku agen agar tetap berada dalam batasan domain ahli patologi tanaman. Agen menunjukkan kemampuan untuk menentukan strategi deteksi yang tepat berdasarkan input pengguna serta menolak permintaan yang tidak relevan atau tidak aman.
    \item Evaluasi sistem menggunakan pendekatan \textit{LLM-as-a-judge} menunjukkan kinerja yang andal pada berbagai skenario pengujian. Pada data distribusi internal (\textit{In-Distribution}), sistem mencapai tingkat keberhasilan tujuan (\textit{Agent Goal Accuracy}) rata-rata sebesar \qty{96}{\percent}. Sistem juga menunjukkan ketahanan yang baik terhadap interaksi multi-tahap dan kasus batas (\textit{edge cases}), dengan skor kepatuhan domain yang sempurna. Meskipun demikian, akurasi identifikasi masih mengalami penurunan pada data di luar distribusi (\textit{Out-of-Distribution}) yang belum terindeks dalam basis pengetahuan.
\end{enumerate}

\section{Saran}
Berdasarkan hasil penelitian dan simpulan yang diperoleh, berikut adalah saran yang dapat diajukan untuk pengembangan sistem lebih lanjut:

\begin{enumerate}[nosep]
    \item Perluasan cakupan basis data pengetahuan dan indeks vektor penyakit tanaman sangat diperlukan untuk meningkatkan akurasi pada data \textit{Out-of-Distribution} (OOD). Penambahan variasi data, khususnya untuk penyakit dengan gejala visual yang mirip, akan membantu sistem dalam memberikan diagnosis yang lebih spesifik dan mengurangi ambiguitas.
    \item Optimasi kinerja waktu inferensi pada modul deteksi \textit{open-vocabulary} perlu menjadi perhatian. Penggunaan model OWLv2 saat ini masih memiliki latensi yang cukup tinggi. Eksplorasi teknik kuantisasi model atau penggunaan arsitektur deteksi \textit{open-vocabulary} yang lebih ringan dapat dipertimbangkan untuk meningkatkan responsivitas sistem.
    \item Pengembangan dukungan bahasa yang lebih luas, khususnya Bahasa Indonesia, pada komponen VLM dan basis pengetahuan. Saat ini sistem masih dominan menggunakan Bahasa Inggris karena keterbatasan model \textit{foundation}. Adaptasi ke bahasa lokal akan meningkatkan aksesibilitas dan kebermanfaatan sistem bagi petani di Indonesia.
    \item Pelaksanaan evaluasi secara daring (\textit{online evaluation}) dengan melibatkan pengguna akhir, seperti petani atau penyuluh pertanian, disarankan untuk mengukur efektivitas dan kegunaan sistem dalam kondisi lapangan yang sebenarnya. Umpan balik dari pengguna nyata akan memberikan wawasan berharga yang tidak tertangkap dalam pengujian simulasi.
    \item Pengumpulan dan penyusunan dataset khusus untuk evaluasi \textit{Visual Question Answering} (VQA) multitahap perlu dilakukan. Keterbatasan dataset publik yang umumnya berfokus pada interaksi tunggal menyulitkan pengukuran kemampuan agen dalam mempertahankan konteks dialog visual yang kompleks. Dataset multitahap yang terstandarisasi akan memungkinkan evaluasi yang lebih akurat terhadap kinerja sistem dalam menangani percakapan interaktif yang berkelanjutan.
\end{enumerate}
